\section{The size of a negative excursion in the twisted prime sum}
\label{sec:osc}
The file \texttt{New/Reduction/Zeros/PrimeError.lean} computes
the poles of a smoothed prime sum used in the proof. We ask how far
its normalized error can go below zero when the auxiliary prime varies.
The answer for modulus one has an absolute uncertainty of less than
$0.047$ on either side of an explicit center. It also gives a concrete
question about relations between zeros on the critical line.

Fix a prime $d\equiv7\pmod 8$, let $\chi(n)=(n/d)$, and let
$\gamma_E$ be Euler's constant. With this notation the repository's
\texttt{primeRemainder} at modulus one is
\[
 R_d(x)=\sum_{n\le x}\Lambda(n)(1-\chi(n))(1/n-1/x)
       -\log x+1+\gamma_E-L'(1,\chi)/L(1,\chi).
\]
Put $B_0=2+\gamma_E-\log(4\pi)=0.0461914179\ldots$ and
$B_d=\log(d/\pi)-\gamma_E+2 L'(1,\chi)/L(1,\chi)$.

\paragraph{Excursion theorem.}
Under RH and Dirichlet GRH, for every such $d$,
\[
 \boxed{\quad -B_d-B_0\ \le\ \liminf_{x\to\infty}\sqrt x\,R_d(x)
                   \ \le\ -B_d+B_0.\quad}
\]
Moreover a single sequence $x_j\to\infty$ realizes
$\sqrt{x_j}R_d(x_j)\to B_0-B_d$ simultaneously for every $d$.
In particular the lower limit divided by $\log d$ tends to $-1$ as
$d\to\infty$ through these primes. The two limits here are taken in order,
first in $x$ for fixed $d$. The theorem makes no prediction about the
first $x$ at which a given excursion occurs.

\paragraph{The explicit formula.}
For $F=\zeta$ or $F=L(\cdot,\chi)$ write
\[
 H_F(u)=\sum_{\rho=1/2+i\gamma}
          \frac{e^{i\gamma u}}{1/4+\gamma^2}.
\]
Take all nontrivial zeros with multiplicities, including central zeros.
The zero count $O_d(T\log(2T))$ gives absolute uniform convergence.
Complex conjugation makes each series real. For fixed $d$,
\[
 \sqrt x\,R_d(x)=H_\zeta(\log x)-H_{L(\cdot,\chi)}(\log x)
                         +O_d(x^{-1/2}).                         \tag{O1}
\]
Indeed integrate $x^s A(s)/(s(s-1))$ on $\Re s=2$, with
$A=-\zeta'/\zeta+L'/L(\cdot,\chi)$ and measure $ds/(2\pi i)$.
Mellin inversion gives $\sum_{n\le x}\Lambda(n)(1-\chi(n))(x/n-1)$,
with zero weight at $n=x$. At $s=1$ the residue is
$x(\log x-1-\gamma_E+L'(1,\chi)/L(1,\chi))$. At $\rho$ it is
$x^\rho(m_\zeta(\rho)-m_L(\rho))/(\rho(1-\rho))$, including
coincident zeros. Shift the contour to $\Re s=-1/2$. The kernel
has a simple pole at zero since $\zeta(0)L(0,\chi)\ne0$ for this odd
primitive character. Its residue is a constant depending on $d$.
On the new vertical line the logarithmic derivatives are
$O_d(\log(|\Im s|+2))$ by the functional equations.
Horizontal edges may be sent to infinity along heights separated by
$\gg_d 1/\log T$ from both zero sets in $[T-1,T+2]$, using the
standard local zero count $O_d(\log T)$ and the local partial fraction
formula $F'/F(s)=\sum_{|\Im\rho-\Im s|\le1}(s-\rho)^{-1}
+O_d(\log T)$ in this fixed strip at large heights. Thus the edge
integrals vanish and the new integral is $O_d(x^{-1/2})$. Dividing
the shifted identity by $\sqrt x$ proves (O1).

\paragraph{Mass and recurrence.}
For the completed entire functions $\Xi_F$ (including the factor
$s(s-1)/2$ for zeta), the Hadamard product gives
\[
 \frac{\Xi'_F}{\Xi_F}(1)-\frac{\Xi'_F}{\Xi_F}(0)
    =\sum_{\rho}\frac{1}{\rho(1-\rho)}
    =H_F(0).
\]
The functional equation makes the left side twice the derivative
ratio at $1$. Substitution of the gamma factors gives $B_0$ for zeta
and $B_d$ for the odd quadratic character.
Consequently $|H_F(u)|\le H_F(0)$. This proves the lower bound in the box.
For finitely many real ordinates, simultaneous approximation returns
all their phases arbitrarily close to $1$ at arbitrarily large positive
times, regardless of rational relations. This follows by compactness:
the integer iterates in the finite torus have a convergent subsequence,
and differences of increasingly separated indices return to the identity.
Enumerate all ordinates for all the countably many characters and zeta,
then take $u_j>j$ approximating the first $j$ phases to within $1/j$.
Absolute convergence in each fixed series gives (O1) along $x_j=e^{u_j}$,
which proves the upper bound and the simultaneous assertion. There is
even a single sequence of integer exponents $n_j\to\infty$ such that
$\sqrt{a^{n_j}}R_d(a^{n_j})\to B_0-B_d$ for every integer $a\ge2$
and every eligible $d$: apply the same integer-iterate argument to the
countably many frequencies $\gamma\log a$. In particular the boxed bounds
hold along powers of any given integer base. Finally
the GRH estimate $L'(1,\chi)/L(1,\chi)=O(\log\log d)$ gives the conductor
asymptotic. Ihara, Murty and Shimura prove this estimate and an explicit
formula for this same ramp kernel in
\href{https://www.impan.pl/shop/en/publication/transaction/download/product/82940}
{\emph{On the logarithmic derivatives of Dirichlet $L$-functions at $s=1$}},
Acta Arith.\ 137 (2009), Corollary 3.3.2 and \S4.2. The contribution
here is the two-sided localization and common recurrence for the
difference used by the Lean development.

\paragraph{How much information is lost in the interval?}
The lower endpoint is attained as a limiting value precisely when every
finite integer relation among the positive ordinates of zeta and of
$L(\cdot,\chi)$ has even total coefficient on the zeta ordinates.
Treat each occurrence as labeled by its function, so a shared positive
ordinate already violates the condition. To see the assertion, equality
forces each positive-weight cosine of zeta to approach $-1$ and each
cosine of $L$ to approach $+1$: their individually nonnegative deficits
sum to zero in the limit. Zeta has no central zero, since the alternating
eta series is positive at $1/2$. Central zeros of $L$ already contribute
with the required sign. On each finite set of frequencies, Kronecker's
approximation theorem identifies the closure of the real flow with the
subtorus satisfying its integer relations. The desired sign vector
belongs to that closure exactly under the stated parity condition.
A diagonal sequence and the summable tails give sufficiency for all
ordinates. This identifies the additional arithmetic input that would
give an exact lower endpoint. Whether the parity condition holds for
these spectra is unresolved here.

\paragraph{Numerical examples.}
For $d=7$ the endpoints are approximately $-0.30142740$ and
$-0.20904456$. For $d=23$ they are $-0.61549819$ and $-0.52311535$.
These values evaluate the proven bounds; they do not estimate a lower
limit by finite sampling. The companion script uses the finite
Hurwitz-zeta formula at zero and the functional equation. In detail, put
$A=-d^{-1}\sum_{a=1}^{d-1}a\chi(a)=L(0,\chi)$. Then
\[
 \frac{L'(1,\chi)}{L(1,\chi)}
 =\log(2\pi)+\gamma_E
  -\frac{1}{A}\sum_{a=1}^{d-1}\chi(a)\log\Gamma(a/d).
\]
